\section{Implementation of the Core Subsystems}
\label{sec:core_implementation}

Once the architectural model has been defined, the next question is how it is realised in terms of concrete C++ constructs and how the individual subsystems interact. This chapter examines the principal building blocks of the library and follows them from entity descriptions to query execution, prepared queries, migrations, and the enforcement of constraints.

The focus here is no longer on layers as an abstract architectural scheme, but on their actual mechanics. How exactly are \code{where} clauses accumulated? How is a query object preserved for repeated use? How does \code{migrate<DB>} transform the difference between the entity model and \code{live\_schema} into DDL operations? How does the thread-safety layer guarantee correct use of connections? The answers to these questions show how far the library is developed as a real engineering system.

\subsection{Repository-backed mini case study: \code{User} and \code{Post}}

To keep the discussion anchored in the real repository, this chapter uses as its main mini case study the \code{User} and \code{Post} types from \path{tests/integration/test_mockdb.cpp}. They are especially useful because they combine scalar fields, table naming, and a join scenario. The same example is later reused for prepared queries, result typing, and SQL rendering assertions.

\begin{lstlisting}[language=C++,caption={Entity types from \protect\path{tests/integration/test_mockdb.cpp}},label={lst:mockdb_entities}]
struct Post
{
    orm::property<int, "id">             id;
    orm::property<int, "author_id">      author_id;
    orm::property<std::u8string, "body"> body;
};

struct User
{
    orm::property<int, "id">            id;
    orm::property<std::u8string, "name"> name;
    orm::property<double, "score">      score;
};
\end{lstlisting}

Listing~\ref{lst:mockdb_entities} shows not only how the entity model is described, but also how it participates directly in the tests. \code{table\_name\_trait<User>} and \code{table\_name\_trait<Post>} fix the logical containers \code{"users"} and \code{"posts"}, while member pointers such as \code{\&User::id} and \code{\&Post::author\_id} later appear in the query builder. This is a good illustration of one of the central themes of the thesis: the model, the query form, and the execution tests all rely on the same compile-time artefacts.

\subsection{Scalar properties and naming of storage containers}

The implementation of \code{property<T, Name>} binds a value type to a compile-time name. The \code{column\_name()} function provides unambiguous access to the symbolic field name and is used by connectors when rendering toward SQL columns, document fields, or other backend-specific notions. This construction is both minimalistic and sufficient because it does not introduce an external metadata registry.

Naming of the target storage container is realised through \code{table\_name\_trait<T>}. Although the name is inherited from ORM tradition, architecturally it plays a more general role: in SQLite and PostgreSQL it denotes a table, in MongoDB a collection, in Redis a key prefix, and in Neo4j a label. In this way, the library uses a unified logical interface for addressing physically distinct structures.

In practical terms, this means that a connector author does not extract information from an external schema or from a runtime registry, but from the type system itself. This reduces boilerplate and makes naming errors easier to localise during compilation or in unit tests.

\subsection{Field wrappers, rules, and placeholders}

A field in a query expression is represented by \code{field<\&T::m>}. This is a \code{constexpr} wrapper over a member pointer that carries enough information about the table type, the value type, and the column name. Operators over \code{field} construct typed \code{Rule} objects. Expressions such as \code{field<\&User::id> == Placeholder<int>\{\}} are therefore not strings, but compile-time descriptions of predicates.

The placeholder mechanism has two forms. \code{Placeholder<T>} models anonymous parameters that are consumed positionally from \code{execute()} arguments. \code{orm::ph<T, std::placeholders::\_N>} models indexed parameters that allow a single runtime value to be reused in more than one position. This is especially important for backends such as SQLite and PostgreSQL, where indexed placeholders are naturally supported.

The tests \code{IndexedPlaceholderRendersQuestionMarkN}, \code{IndexedPlaceholderReusedSameArgInSQL}, and \code{TwoDistinctIndexedPlaceholdersRenderCorrectSlots} in \path{tests/integration/test_mockdb.cpp} show that this subsystem is not merely theoretical ornamentation, but a real mechanism for expressing more precise parameter-binding scenarios.

\subsection{Construction of query objects}

The factory functions \code{select}, \code{insert}, \code{update}, and \code{deleteq} create query objects onto which \code{where}, \code{join}, \code{order\_by}, \code{group\_by}, and \code{limit} are layered successively. Crucially, each of these operations does not modify a string, but returns a new typed object with extended internal state. As a result, the query builder is fluent in appearance and static in nature.

This allows the tests to check not only the final rendering result, but the very shape of the intermediate representation --- for example, the presence of an \code{INNER JOIN}, the order clause, or the resulting response tuple type. In \path{tests/integration/test_mockdb.cpp}, this appears in \code{SelectWithInnerJoin}, \code{SelectWithOrderByDesc}, \code{SelectWithLimit}, and many other scenarios that validate different phases of query composition.

\begin{figure}[H]
\centering
\begin{tikzpicture}[node distance=1.35cm and 1.2cm]
    \node[ormaccent] (api) {\textbf{Fluent API}\\\code{select}/\code{where}/\code{join}/\code{limit}};
    \node[ormblock, below=of api] (typed) {New typed\\query state at each step};
    \node[ormblock, below left=of typed] (shape) {Compile-time\\query structure};
    \node[ormblock, below right=of typed] (params) {Runtime\\parameter values};
    \node[ormblock, below=of $(shape)!0.5!(params)$] (exec) {\code{connector\_trait<DB>::execute}};

    \draw[ormline] (api) -- (typed);
    \draw[ormline] (typed) -- (shape);
    \draw[ormline] (typed) -- (params);
    \draw[ormline] (shape) -- (exec);
    \draw[ormline] (params) -- (exec);
\end{tikzpicture}
\caption{Interaction between query composition, compile-time structure, and runtime parameters}
\label{fig:query_composition_execution}
\end{figure}

Figure~\ref{fig:query_composition_execution} shows that the implementation operates in two modes simultaneously: the structural form of the query is fixed in the type, while the concrete values are supplied separately. This is exactly what makes it possible to combine compile-time checking with runtime parameterisation.

\subsection{Execution, results, and field access}

At execution time, \code{db<DB>} forwards the query IR to \code{connector\_trait<DB>::execute}. The result is returned as \code{orm::result<Row, Fields>}, where \code{Row} is a tuple type and \code{fields} preserves information about the selected fields. This design supports both iteration over materialised rows and field access through \code{get\_field<\&T::m>}. The implementation therefore remains type-safe even after the query has been executed.

The tests \code{SelectResultRowTypeIsTuple}, \code{SelectMultiFieldRowType}, \code{SelectGetByMemberPointer}, and \code{SelectGetByIndex} show that the result layer is not merely a container of anonymous rows, but a structured type-safe abstraction. This is especially important for the thesis because it demonstrates that type discipline does not end with the query builder, but continues into result handling.

\subsection{Prepared queries as a reusable execution artefact}

\code{prepared\_query<DB, Query>} extends the execution model. Instead of caching an SQL string, it stores a reference to the backend object together with the query IR itself. This is a natural continuation of the compile-time approach: the library treats the query object as the primary artefact rather than as a temporary step on the way to text.

\begin{lstlisting}[language=C++,caption={Excerpt from \code{prepared\_query<DB, Query>}},label={lst:prepared_query}]
template <typename DB, typename Query>
class prepared_query
{
public:
    prepared_query(DB& conn, Query q) : conn_(conn), query_(std::move(q)) {}

    template <typename... Params>
    auto execute(Params&&... params) const
    {
        return connector_trait<DB>::execute(
            conn_, query_, std::forward<Params>(params)...);
    }

    [[nodiscard]] const Query& query() const noexcept { return query_; }
};
\end{lstlisting}

Listing~\ref{lst:prepared_query} reveals an important engineering choice: reuse is organised around the IR rather than around rendered text. This enables one and the same logical query to be executed repeatedly with different parameter values, without reconstructing the builder form.

\begin{figure}[H]
\centering
\begin{tikzpicture}[node distance=1.35cm and 1.2cm]
    \node[ormaccent] (build) {Query-IR construction};
    \node[ormblock, right=of build] (prepare) {\code{db.prepare(query)}};
    \node[ormblock, right=of prepare] (store) {Stored IR +\\reference to \code{DB}};
    \node[ormblock, below=of store] (exec1) {\code{execute(1)}};
    \node[ormblock, left=of exec1] (exec2) {\code{execute(99)}};
    \node[ormblock, right=of exec1] (exec3) {\code{execute(7)}};

    \draw[ormline] (build) -- (prepare);
    \draw[ormline] (prepare) -- (store);
    \draw[ormline] (store) -- (exec1);
    \draw[ormline] (store) -- (exec2);
    \draw[ormline] (store) -- (exec3);
\end{tikzpicture}
\caption{Lifecycle of \code{prepared\_query}: one IR, multiple executions}
\label{fig:prepared_query_lifecycle}
\end{figure}

This behaviour is precisely what is validated by \code{PrepareReturnsExecutableObject}, \code{PrepareWithParamExecutesCorrectly}, \code{PrepareCanBeCalledMultipleTimesWithDifferentParams}, and \code{PrepareExposesQueryIR} in \path{tests/integration/test_mockdb.cpp}. These tests prove that the prepared-query object is both executable and inspectable, which is strong evidence of internal consistency in the subsystem.

\subsection{Migration layer and DDL generation}

The implementation of \code{migrate<DB>} compares the schema described by entity types against \code{live\_schema}. Whenever a difference exists, DDL operations are produced and delegated to \code{connector\_trait<DB>::ddl\_for(...)}. This is a good example of clean separation of responsibilities: diff logic is shared, while concrete DDL syntax remains backend-specific.

\begin{lstlisting}[language=C++,caption={Excerpt from the migration diff pipeline},label={lst:migration_diff}]
[[nodiscard]] std::vector<ddl_op> diff(
    const live_schema& live,
    std::span<const entity_meta> entities) const
{
    std::vector<ddl_op> ops;

    for (const auto& entity : entities)
    {
        auto live_it = live.find(entity.table_name);
        if (live_it == live.end())
        {
            ops.push_back(create_table_op{entity.table_name, entity.columns});
            continue;
        }

        for (const auto& [col, type] : entity.columns)
        {
            if (live_it->second.find(col) == live_it->second.end())
                ops.push_back(add_column_op{entity.table_name, col, type});
        }
    }
    return ops;
}
\end{lstlisting}

It is precisely here that the difference between a \emph{minimal working connector} and a \emph{complete connector} becomes visible. A backend may execute queries without supporting DDL generation. But if it is to participate in schema-aware tooling, its contract must be extended.

\begin{figure}[H]
\centering
\begin{tikzpicture}[node distance=1.35cm and 1.2cm]
    \node[ormaccent] (entity) {Entity metadata\\\code{entity\_meta}};
    \node[ormblock, right=of entity] (live) {Runtime schema\\\code{live\_schema}};
    \node[ormblock, below=of $(entity)!0.5!(live)$] (diff) {\code{migrate<DB>::diff}};
    \node[ormblock, below left=of diff] (ops) {DDL operations\\\code{create/add/drop/alter}};
    \node[ormblock, below right=of diff] (ddl) {\code{generate\_ddl}};
    \node[ormblock, below=of $(ops)!0.5!(ddl)$] (run) {dry-run or execution path};

    \draw[ormline] (entity) -- (diff);
    \draw[ormline] (live) -- (diff);
    \draw[ormline] (diff) -- (ops);
    \draw[ormline] (diff) -- (ddl);
    \draw[ormline] (ops) -- (run);
    \draw[ormline] (ddl) -- (run);
\end{tikzpicture}
\caption{Flow of the schema-migration subsystem}
\label{fig:migration_pipeline}
\end{figure}

The tests in \path{tests/unit/test_schema_migration.cpp} cover table creation, addition and removal of columns, dry-run exit codes, DDL generation, and the state-machine behaviour of \code{run()}. This is a good example of how higher-responsibility library subsystems are validated through targeted unit tests rather than through integration smoke tests alone.

\subsection{Thread safety and transaction helpers}

The library also contains a dedicated layer for \code{connection\_pool<DB, N>}, \code{thread\_local\_db<DB>}, and \code{transaction\_guard<DB>}. These components show that the implementation does not stop at one-off query execution, but addresses more realistic operational scenarios. The important architectural decision is that access to these mechanisms is also capability-based. For example, \code{connection\_pool} requires \code{supports\_concurrent\_execute}, while \code{transaction\_guard} requires \code{supports\_transactions}.

\begin{lstlisting}[language=C++,caption={Excerpt from the thread-safety helper layer},label={lst:thread_safety_helpers}]
template <typename DB, std::size_t N>
class connection_pool
{
    static_assert(
     requires { typename connector_trait<DB>::supports_concurrent_execute; },
     "connection_pool requires connector_trait<DB>::supports_concurrent_execute.");

public:
    [[nodiscard]] connection_guard<DB> acquire();
};

template <typename DB>
class transaction_guard
{
    static_assert(
     requires { typename connector_trait<DB>::supports_transactions; },
     "transaction_guard requires connector_trait<DB>::supports_transactions.");
};
\end{lstlisting}

The compile-time contract therefore does not end at query form, but extends into the lifecycle policies of the connector. This matters for the thesis claim that the library is not merely a query builder, but an infrastructure for systematic and safe backend interaction.

\begin{table}[H]
\centering
\small
\caption{Core subsystems and their test evidence}
\label{tab:core_subsystem_tests}
\begin{tabularx}{\textwidth}{L{3.1cm}L{3.8cm}Y}
\toprule
\textbf{Subsystem} & \textbf{Main files} & \textbf{Test evidence} \\
\midrule
Entity modelling & \path{ORM/entity/property.hpp}, \path{ORM/entity/relationship.hpp} & \path{tests/unit/test_property.cpp}, \path{tests/unit/test_relationship.cpp} \\
Query composition & \path{ORM/query/select.hpp} and related query headers & \path{tests/unit/test_query.cpp}, \path{tests/integration/test_mockdb.cpp} \\
Prepared queries & \path{ORM/connector/prepared_query.hpp} & the \code{Prepare*} scenarios in \path{tests/integration/test_mockdb.cpp} \\
Schema migration & \path{ORM/db/migration/migration.hpp} & \path{tests/unit/test_schema_migration.cpp} \\
Thread safety & \path{ThreadSafety/thread_safety.hpp} & \path{tests/unit/test_thread_safety.cpp} \\
\bottomrule
\end{tabularx}
\end{table}

\subsection{Implementation conclusion}

The core subsystems have direct support in the code and in the tests. \path{tests/unit/test_property.cpp} and \path{tests/unit/test_relationship.cpp} validate the entity descriptions. \path{tests/unit/test_query.cpp} demonstrates that the query builder does indeed construct the intended compile-time structures. \path{tests/integration/test_mockdb.cpp} and \path{tests/integration/test_sqlite.cpp} validate execution and prepared-query scenarios. \path{tests/unit/test_schema_migration.cpp} covers the migration diff logic, while \path{tests/unit/test_thread_safety.cpp} verifies the concurrency layer.

Consequently, the implementation of the core subsystems is not merely a project intention, but an architectural solution supported by systematic verification. This is the most important conclusion of the present chapter: the compile-time ORM architecture in the repository does not remain at the level of abstraction, but materialises into working subsystems with clear responsibilities and explicit test coverage.
